\section{Conclusion}
\label{sec:conclusion}

We have presented \texttt{micrograd}, an open-source FEniCSx framework
for density-based topology optimisation of coupled
Brinkman--convection-diffusion systems.
The implementation makes four documented contributions:

\begin{enumerate}
    \item \textbf{Continuous adjoint implementation.}
          A complete FEniCSx implementation of the continuous adjoint
          for coupled Brinkman--convection-diffusion, following
          standard Lagrangian methodology~\cite{borrvall2003},
          with the momentum--concentration coupling $-\lambda\nabla c$
          and misfit-derived Dirichlet outlet condition.

    \item \textbf{SUPG stabilisation of forward and adjoint equations.}
          Standard SUPG~\cite{brooks1982} applied consistently
          to both equations in FEniCSx UFL, reducing sensitivity-field
          oscillations from $\sim0.4$ to below $10^{-3}$
          at $\Pe\sim10^2$--$10^5$.

    \item \textbf{Documented continuation schedule.}
          Five components (Helmholtz filtering, Heaviside
          $\beta$-sharpening, $\alpha_{\max}$ ramping, hybrid OC/MMA,
          sinusoidal initialisation) with each component's
          failure mode documented (Table~\ref{tab:continuation});
          reduces gray fraction from $\grayShortRun$ to $\grayFrac$.

    \item \textbf{Open-source reproducible release.}
          MIT licence, Docker, CI, and Zenodo archival;
          all results reproducible from a single command
          in approximately one hour on a standard laptop.
\end{enumerate}

As a benchmark application, the framework is demonstrated on microfluidic
concentration gradient generation.
On an $80\times20$ mesh, the Brinkman surrogate optimum achieves
RMSE~$=\rmseTopOpt=0.058$ for a linear target profile, with gray-zone
fraction $\grayFrac=0.074$.
Threshold-projection to a binary-permeability field increases RMSE
from $\rmseTopOpt$ to $\rmseBinaryBrinkman$.
The optimised design exploits the continuous permeability field of the
Brinkman model to achieve precise mixing targets;
this porous optimum is not representable as a binary solid/fluid
geometry without performance penalty.
\texttt{micrograd} is deliberately a surrogate tool for rapid
inverse design exploration in porous media: it identifies
high-performance continuous-permeability concepts that inform
subsequent binary design or experimental iteration.
This two-stage workflow is standard in Brinkman topology
optimisation~\cite{borrvall2003} and is the intended use case.

\paragraph{Future work.}
The principal open directions are:
(i)~binary realisation via robust projection~\cite{wang2011} or
iterative thresholding;
(ii)~a discrete adjoint for exact gradient magnitude;
(iii)~three-dimensional extension and experimental fabrication;
(iv)~adaptive mesh refinement targeting the diffusion sublayer
($\delta\approx 4\,\mu$m).

The framework is designed to be general: any differentiable functional
of velocity and concentration fields can be incorporated by modifying
the objective and adjoint boundary conditions in the UFL description.
This generality, combined with a computational cost of under one hour on a standard laptop
(or under thirty minutes for the 600-iteration quick run)
and full end-to-end reproducibility,
makes \texttt{micrograd} a practical foundation for surrogate-based
computational microfluidic optimisation and a reference implementation
for coupled flow-transport adjoint methods in FEniCSx.
